\chapter*{Prefácio}
\addcontentsline{toc}{chapter}{Prefácio}

\hayashi{} nasceu de uma insatisfação concreta: as ferramentas existentes para econometria
aplicada exigem que o pesquisador escolha entre expressividade e desempenho, entre uma
sintaxe próxima da notação estatística e a capacidade de automatizar análises complexas.

Stata tem a sintaxe certa, mas é proprietário, fechado e não é uma linguagem de propósito
geral. R tem o ecossistema certo, mas a sintaxe é inconsistente e o modelo de dados é
surpreendente para quem vem de Stata. Python tem a flexibilidade certa, mas os pacotes
estatísticos tratam econometria como subcaso da aprendizagem de máquina.

\hayashi{} é uma linguagem interpretada, de propósito geral, construída especificamente
para análise econométrica. A sintaxe homenageia Stata nos lugares certos e adota idiomas
modernos onde Stata mostra sua idade. O motor é escrito em Rust — sem máquina virtual,
sem \textit{garbage collector}, sem tempo de inicialização perceptível.

\bigskip

Este livro documenta tanto a linguagem quanto o núcleo econométrico: da estrutura de
tipos e controle de fluxo aos estimadores de séries temporais, inferência causal e
métodos avançados. É um retrato do estado atual do projeto — um documento vivo, que
evolui com o software.

\bigskip
\noindent\rule{\textwidth}{0.4pt}
\bigskip

\section*{Um projeto de um desenvolvedor — e de toda a comunidade}

\hayashi{} é, até o momento, obra de um único desenvolvedor. Isso tem consequências
práticas que o leitor deve conhecer.

Do lado positivo, há coerência de design: cada decisão de sintaxe, cada interface de
estimador, cada mensagem de erro foi tomada com o mesmo conjunto de valores em mente.
Não há comitês, não há \textit{bikeshedding}, não há camadas de compatibilidade
acumuladas por décadas. O código é enxuto porque ninguém adicionou nada ``só para não
perder''.

Do lado da limitação, um par de olhos — por mais criterioso que seja — não vê tudo.
O projeto conta com três camadas de verificação: testes unitários cobrindo os
primitivos da linguagem; scripts de integração que exercitam a \textsc{Hayashi} como
caixa preta; e scripts \textit{DGP} (processo gerador de dados) para cada estimador,
onde o parâmetro verdadeiro é conhecido e a recuperação pode ser verificada
numericamente. Ainda assim, casos de borda existem, e alguns certamente escaparam.

\bigskip

\noindent\textbf{Licença.} \hayashi{} é software livre, distribuído sob os termos da
\textbf{GNU General Public License versão~3} (GPL~v3). Isso significa que você pode
usar, estudar, modificar e redistribuir o programa — inclusive para fins comerciais —
desde que qualquer trabalho derivado seja igualmente livre. O código-fonte é público
e pertence, no sentido mais literal, a quem o usa.

A escolha da GPL~v3 não é acidental. Ferramentas econométricas proprietárias criam
dependência institucional: universidades e laboratórios ficam reféns de licenças, de
políticas de preço e de decisões corporativas sobre o que o software pode ou não fazer.
Software livre quebra esse ciclo. \hayashi{} deve ser acessível ao pesquisador em
qualquer lugar do mundo, independentemente de financiamento.

\bigskip

\noindent\textbf{Colaboradores e testadores são bem-vindos.} Se você encontrou um
resultado inesperado, um estimador que não converge quando deveria, uma mensagem de
erro obscura ou simplesmente uma sintaxe que poderia ser mais clara — isso é
informação valiosa. Abra uma \textit{issue} no repositório. Se quiser ir além, uma
correção bem testada ou um novo capítulo de documentação é contribuição tão legítima
quanto qualquer linha de código Rust.

O projeto precisa especialmente de pessoas dispostas a testar \hayashi{} em dados
reais: séries de tempo com sazonalidade, painéis desbalanceados, modelos com
multicolinearidade, dados com missings. Esses são os casos que os scripts DGP não
cobrem, porque dados sintéticos são, por construção, bem-comportados.

\hayashi{} só se torna robusto com uso. E uso honesto — que inclui reportar o que
quebra — é a forma mais direta de colaborar.

\bigskip

\noindent\textit{Junho de 2026.}
